% Generated draft: requires mathematical and source review.
\documentclass[11pt]{article}
\usepackage[margin=25mm]{geometry}
\usepackage{fontspec}
\setmainfont{Latin Modern Roman}
\usepackage{amsmath,amssymb,mathtools,unicode-math}
\setmathfont{Latin Modern Math}
\usepackage{enumitem,needspace,xurl,hyperref}
\setlength{\parindent}{0pt}
\setlength{\parskip}{5pt}
\setlength{\emergencystretch}{3em}
\newcommand{\N}{\mathbb{N}}
\newcommand{\Z}{\mathbb{Z}}
\newcommand{\Q}{\mathbb{Q}}
\newcommand{\R}{\mathbb{R}}
\newcommand{\C}{\mathbb{C}}
\newcommand{\F}{\mathbb{F}}
\newcommand{\A}{\mathbb{A}}
\newcommand{\PP}{\mathbb{P}}
\newcommand{\E}{\mathbb{E}}
\newcommand{\eps}{\varepsilon}
\newcommand{\dd}{\,\mathrm{d}}
\begin{document}
\section*{UnsolvedMath ID: 6900005}
\textbf{Source problem code:} catalogue entry
\subsection{Definitions and mathematical statement}
For context, a tensegrity framework can be described by a graph $G=(V,E)$ and vertex positions $p_v\in\R^d$; an equilibrium self-stress is a choice of edge scalars $\omega_{uv}$ satisfying
\[
 \sum_{v:uv\in E}\omega_{uv}(p_v-p_u)=0
 \qquad(u\in V).
\]
The source's configuration-space stratification may retain more than the dimension of this stress space, such as sign or equivalence data; its actual definition is not included in the supplied record. No stress-rank stratification is silently substituted.

Let $B_G$ denote the intended configuration space and $\mathcal S_G$ its source-defined strata. The task is to identify all strata of codimension greater than one that are not defined as intersections of closures of several codimension-one strata. Under a literal set-equality reading, one would seek those $S\in\mathcal S_G$ with $\operatorname{codim}_{B_G}S\ge2$ for which no finite family of codimension-one strata $T_1,\ldots,T_r$ satisfies
\[
 S=\overline{T_1}\cap\cdots\cap\overline{T_r}.
\]
However such an intersection is closed while a stratum need not be. The intended criterion may concern the closure of $S$, a component, or a local defining intersection rather than literal equality. The record does not choose among these, so the displayed equality explains only the literal reading for review.

The graph range, ambient dimension, permitted coincident vertices, quotienting by Euclidean or affine transformations, codimension convention, closure topology, and exact stratification must be supplied before a complete characterization is determined. The research summary gives a two-vertex single-edge framework as an example with a codimension-two stratum, but does not specify its ambient or quotient conventions.
\subsection{Short English statement}
Which higher-codimension tensegrity strata are genuinely new rather than intersections of closures of codimension-one strata? The intersection criterion and the underlying stratification need their exact source definitions.

The supplied summary reports known exceptional examples, including the two-vertex single-edge case, but no complete characterization. It labels the question open. The OPEN-TRIAGE assessment and 2020 author-manuscript problem are retained with current status marked for review.
\subsection{Sources}
\begin{itemize}
\item Local UnsolvedMath registry, permanent ID 6900005, source code catalogue entry, ``Configuration Spaces of Tensegrities --- Problem 5,'' proposed by Oleg Karpenkov, 2020. Dataset curation and metadata: Ulam AI and UnsolvedMath Contributors, CC BY 4.0.
\item Oleg Karpenkov, \emph{Open Problems on Configuration Spaces of Tensegrities} (2020), Problem 5, PDF page 4, \url{https://livrepository.liverpool.ac.uk/3022838/1/stratum.pdf}; normalized pdf-text extraction, accessed 2026-07-29.
\item The supplied summary gives the single-edge example and an open-status assessment without a later primary resolution reference. Source-specific provenance: the public institutional-repository author manuscript labels the item open; current status and reuse terms NEEDS\_REVIEW, separately from the dataset licence.
\end{itemize}
\end{document}
